\section{Training}
\label{sec:training}

This section describes the training protocol and reports the training record. Every number in it
is a generated macro; a value that the training record does not contain yet is printed as
\NotAvailable{} rather than estimated.

\paragraph{Data and batching.} Training windows are drawn from the training blocks of the spatial
split of Section~\ref{sec:data}, with each batch balanced so that each body contributes a third of
its windows. The training split holds \DatasetMoonWindows{} lunar, \DatasetMarsWindows{} Martian and \DatasetEarthWindows{} terrestrial windows (Section~\ref{sec:data}); the manifest records an upper estimate of about \DatasetCpuHours{} CPU hours for building the dataset. Physics targets are computed for every window by the same fingerprint code that ranks
candidates (Section~\ref{sec:fingerprint}) and standardised with statistics of the training blocks
only. Morphology labels, the crater densities of the lunar and Martian crater catalogues \citep{robbins2019,robbins2012}, are reserved for the probes and never used in training. These labels were not built for this study, so the morphology probe is not available.

\paragraph{Hardware and schedule.} The runs evaluated on test blocks were trained on a CPU server at a reduced scope, without a GPU. Table~\ref{tab:training} gives the schedule as the training record states it; Section~\ref{sec:benchmarks} adds later runs on a laptop accelerator, evaluated on validation blocks only.

\begin{table}[tbp]
\centering
\caption{Training schedule of the first main run, from its training record. The planned schedule was \PlannedEpochs{} epochs at a batch of \PlannedBatch{} in brain floating point with a stem stride of \PlannedStemStride; the run used a stem stride of \UsedStemStride{} in single precision, and its record lists each deviation.}
\label{tab:training}
\footnotesize
\begin{tabular}{@{}ll@{}}
\toprule
Item & Value \\
\midrule
Hardware & CPU only, \ThruCpuCount{} virtual CPUs \\
CPU threads per run & \TrainThreads{} \\
Epochs & \TrainEpochs{} \\
Windows per epoch & \TrainWindowsPerEpoch{} \\
Training rows & \TrainRows{} \\
Batch size & \TrainBatch{} \\
Wall clock hours, first main run & \TrainHours{} \\
CPU core hours, first main run & \TrainCoreHours{} \\
Backbone parameters & \EncoderParams{} \\
\bottomrule
\end{tabular}
\end{table} The
planned grid covers seeds \ConfigSeeds{} of the main configuration and the adversarial weights
\ConfigLambdas; Section~\ref{sec:results-variations} lists which runs exist.

\paragraph{Optimiser.} We use AdamW \citep{loshchilov2017} with a learning rate scaled linearly
with the batch size after a warm up \citep{goyal2017}, cosine decay of the learning rate in the manner of \citet{loshchilov2016} without restarts, and
decoupled weight decay. The adversarial weight follows the schedule of Equation~\eqref{eq:total}.
Model selection uses validation blocks only; the test blocks are read once, after the
configuration is frozen, by a logged command.

\paragraph{Augmentation and regularisation.} The views of Section~\ref{sec:architecture} are the
only augmentation. They contain no amplitude or scale change, because the physics targets must be
identical for both views of a window. VICReg's variance and covariance terms regularise the
embedding against collapse \citep{bardes2021}, and the masked physics loss acts as a supervised
regulariser that ties the embedding to the fingerprint.
